% !PRINTED 44
\weberpage{44}

%<<CONT
darstellbar, so dass
\weq{3}{0 \leqq \alpha_1 < a_1,\ 0 \leqq \alpha_2 < a_2, \ldots 0 \leqq \alpha_\nu < a_\nu.}

Die Exponenten $\alpha_1, \alpha_2 \ldots \alpha_\nu$ oder irgend welche ihnen nach
den Moduln $a_1, a_2 \ldots a_\nu$ congruente Zahlen heissen die \gesperrt{Indices}
des Elementes $A$, und es gilt der Satz:

1. \gesperrt{Man erhält die Indices eines Compositums $A A'$,
wenn man die entsprechenden Indices der beiden
Factoren addirt.}

Wenn wir jedem Elemente $A$ irgendwie einen \gesperrt{Zahlenwerth}
 zuordnen, so können wir diese Zuordnung eine Function
von $A$ nennen. Eine solche Function $\chi(A)$ soll nun ein
\gesperrt{Gruppencharakter} genannt werden, wenn sie für je zwei Elemente
 $A$, $A'$ von $S$ der Bedingung genügt:
\weq{4}{\chi(AA') = \chi(A)\chi(A'),}
und folglich auch der allgemeineren
\[ \chi(AA'A''\ldots) = \chi(A)\chi(A')\chi(A'')\ldots \]
Zwei solche Functionen $\chi$ und $\chi_1$ werden als verschieden
angesehen, wenn es wenigstens ein Element $A$ giebt, wofür $\chi(A)$
von $\chi_1(A)$ verschieden ist. Diese Charaktere wollen wir nun
näher bestimmen.

Setzen wir $A' = 1$, so ergiebt sich aus (4):
\weq{5}{\chi(A) = \chi(A)\chi(1),}
und wenn wir ein für allemal den Fall ausschliessen, dass alle
$\chi(A) = 0$ sind:
\weq{6}{\chi(1) = 1,}
also der erste Satz:

2. \gesperrt{Jeder Charakter hat für das Einheitselement den
Werth 1.}

Setzen wir ferner $A' = A$, so folgt aus (4):
\[ \chi(A^2) = [\chi(A)]^2, \]
und daraus durch vollständige Induction für jeden Exponenten $h$:
\weq{7}{\chi(A^h) = \chi(A)^h.}
Bezeichnen wir also mit $a$ den Grad des Elementes $A$, so
folgt, wenn man in (7) $h = a$ setzt und (6) benutzt:
\weq{8}{\chi(A)^a = 1,}
also:
%CONT>>
